% Experimental LaTEI running-title mapping.
% This file is a compilation artifact, not a reversible source.
\lateiDeclareRunningTitle{Aspects ludiques dans l’appareil héraldique des manuscrits de Léon X (1513-1521)}{Aspects ludiques dans l’appareil héraldique…}
\lateiDeclareRunningTitle{L’héraldique de Jules III Ciocchi del Monte (1550-1555) dans l’ornement pour le livre}{L’héraldique de Jules III Ciocchi del Monte (1550-1555)…}
\lateiDeclareRunningTitle{Érudits et héraldique dans la Rome d’Urbain VIII (1623-1644)}{Érudits et héraldique dans la Rome d’Urbain VIII…}
\lateiDeclareRunningTitle{Francesco Gualdi et les armoiries, entre érudition et conservation}{Francesco Gualdi et les armoiries, entre érudition…}
\lateiDeclareRunningTitle{Les lettrés du milieu romain des Pamphili et l’origine de leur emblème nobiliaire}{Les lettrés du milieu romain des Pamphili et l’origine…}
\lateiDeclareRunningTitle{La « colombe portant le rameau d’olivier au bec » à Gubbio et à Rome}{La « colombe portant le rameau d’olivier au bec »…}
\lateiDeclareRunningTitle{L’écu des Pamphili durant le pontificat d’Innocent X (1644-1655) : deux exemples en architecture}{L’écu des Pamphili durant le pontificat d’Innocent X…}
\lateiDeclareRunningTitle{Épilogue. Une nouvelle fortune pour le blason pamphili : la Raccolta Di Varie Targhe di Roma de Filippo Juvarra (1711)}{Épilogue. Une nouvelle fortune pour le blason pamphili…}
\lateiDeclareRunningTitle{Girolamo Farnèse, légat à Bologne : « Aeterna hic LILIA semper erunt »}{Girolamo Farnèse, légat à Bologne : « Aeterna hic LILIA…}
\lateiDeclareRunningTitle{« Micet in vertice » : Pietro Vidoni et la célébration d’Alexandre VII Chigi}{« Micet in vertice » : Pietro Vidoni et la célébration…}
\lateiDeclareRunningTitle{Piranèse, l’héraldique et Clément XIII Rezzonico (1758-1769)}{Piranèse, l’héraldique et Clément XIII Rezzonico…}
\lateiDeclareRunningTitle{La mise en signe de l’espace par les papes : pratiques et conflits}{La mise en signe de l’espace par les papes : pratiques…}
\lateiDeclareRunningTitle{La collégiale Saint-Jean-Baptiste puis Saint-Louis de Castelnau-Bretenoux : un lien étroit et discret avec la papauté}{La collégiale Saint-Jean-Baptiste puis Saint-Louis…}
\lateiDeclareRunningTitle{Autres témoignages héraldiques intéressants présents dans l’église}{Autres témoignages héraldiques intéressants présents…}
\lateiDeclareRunningTitle{La basilique : un monument marqué par ses constructeurs et ses utilisateurs}{La basilique : un monument marqué par ses constructeurs…}
\lateiDeclareRunningTitle{La basilique lieu du cérémonial : cérémonies éphémères et représentations}{La basilique lieu du cérémonial : cérémonies éphémères…}
\lateiDeclareRunningTitle{L’exercice du pouvoir : quelques exemples d’utilisation des emblèmes}{L’exercice du pouvoir : quelques exemples d’utilisation…}
\lateiDeclareRunningTitle{Dématérialisation, inchoativité, contextualisation : les Apothéoses Barberini (1632-1636) et Médici (1644-1646)}{Dématérialisation, inchoativité, contextualisation…}
\lateiDeclareRunningTitle{L’écu familial et l’ancêtre fondateur : l’apothéose Altieri ( circa 1675-1676)}{L’écu familial et l’ancêtre fondateur : l’apothéose…}
\lateiDeclareRunningTitle{Extension scénique, dissémination et transformisme : l’apothéose Aldobrandini (1596-1602)}{Extension scénique, dissémination et transformisme…}
\lateiDeclareRunningTitle{De la scène au milieu, de l’interprétation à la subjectivation}{De la scène au milieu, de l’interprétation…}
